% concatenated from final_main.tex
\documentclass{amsart}
\usepackage[english]{babel}
\usepackage[utf8]{inputenc}
%\usepackage[T1]{fontenc}
\usepackage{amsmath, amssymb,epic,graphicx,mathrsfs,enumerate}
\usepackage[all]{xy}
\usepackage{xcolor}
\usepackage{comment}
\usepackage{enumitem}
\usepackage{hyperref}
\usepackage[noabbrev,capitalize]{cleveref}
\usepackage[]{frontespizio}
\usepackage{amssymb}
\usepackage{amsmath}
\usepackage{amsthm}
\usepackage{graphicx}
\usepackage{hyperref}
\usepackage{amscd}
\usepackage{enumerate}
\usepackage{mathrsfs}
%\usepackage[notref,notcite]{showkeys}

\renewcommand{\thesection}{\arabic{section}}

%\setlength{\parskip}{2mm}
\usepackage{amsmath}
\usepackage{amsthm}
\usepackage{amssymb}
\usepackage{latexsym}
\usepackage{epsfig}
\DeclareMathOperator{\nil}{\mathcal N}
\DeclareMathOperator{\perm}{Sym}
\DeclareMathOperator{\alt}{Alt}
\DeclareMathOperator{\core}{Core} 
\DeclareMathOperator{\frat}{Frat} 
\DeclareMathOperator{\fit}{Fit} 
\DeclareMathOperator{\aut}{Aut} 
\DeclareMathOperator{\inn}{Inn} 
\DeclareMathOperator{\out}{Out} 
\DeclareMathOperator{\soc}{soc}
\DeclareMathOperator{\dist}{dist}
%\DeclareMathOperator{\nil}{nil}
\DeclareMathOperator{\sol}{sol}
\DeclareMathOperator{\F}{F}
\DeclareMathOperator{\oo}{O}
\newcommand{\op}{\oo_p}
\DeclareMathOperator{\diam}{diam}
\DeclareMathOperator{\kndo}{End}
\DeclareMathOperator{\gl}{GL} 
\DeclareMathOperator{\ssl}{SL} 
\DeclareMathOperator{\ssp}{Sp}
\DeclareMathOperator{\csp}{CSp}
\DeclareMathOperator{\psp}{PSp}
\DeclareMathOperator{\pcsp}{PCSp}
\DeclareMathOperator{\Z}{Z}
\DeclareMathOperator{\syl}{Syl}
\DeclareMathOperator{\B}{B}
\DeclareMathOperator{\I}{I}
\DeclareMathOperator{\R}{R}
\DeclareMathOperator{\C}{C}
\DeclareMathOperator{\nor}{N}
\DeclareMathOperator{\so}{SO}
\DeclareMathOperator{\pgl}{PGL}
\DeclareMathOperator{\pgammal}{P\Gamma L}
\newcommand{\sym}{\mathrm{Sym}}
\newcommand{\supp}{\mathrm{supp}}
\newcommand{\Aut}{\mathrm{Aut}}
\newcommand{\psl}{\mathrm{PSL}}

\DeclareMathOperator{\End}{End}

\newcommand{\gen}[1]{\left\langle#1\right\rangle} 
\newcommand{\al}{\alpha}
\newcommand{\be}{\beta} 
\newcommand{\ol}{\overline}
\newcommand{\N}{\mathbb N}

\newcommand{\st}{such that }
\newcommand{\ifa}{if and only if } 
\newcommand{\sg}{subgroup }
\newcommand{\sgs}{subgroups }
\newcommand{\gr}{group }
\newcommand{\grs}{groups }
\newcommand{\fm}{formation }
\newcommand{\fms}{formations }
\newcommand{\Wh}{We have }


\newcommand{\ov}[1]{\overline{#1}}
\newcommand{\w}[1]{\widetilde{#1}}
\newcommand{\FF}{\mathbb{F}}
\newcommand{\n}{\mathfrak{N}}
\newcommand{\fff}{\mathfrak{F}_2}
\newcommand{\aaa}{\mathfrak{A}}
\newcommand{\ff}{\mathfrak{F}}
\newcommand{\fr}{\phi}
\newcommand{\G}{\Gamma}
\newcommand{\gff}{G^{\ff}}
\newcommand{\ca}{\mathcal}

\newcommand{\la}{\langle}
\newcommand{\ra}{\rangle}
\newcommand{\nn}{\mathrel{\unlhd}}
\newcommand{\sd}{\rtimes}
\newcommand{\x}{\times}
\newcommand{\gf}{\G_\ff}
\newcommand{\iso}{\mathcal{I}_\ff}
\newcommand{\isoo}{\mathcal{I}_{\fff}}
\newcommand{\is}[1]{\mathcal{I}_{#1}}
\newcommand{\cF}{\F}
\newcommand{\rf}{\R_\ff}
\newcommand{\gs}{\G_\mathfrak{S}}
\newcommand{\g}[1]{\G_\mathfrak{#1}}
\newcommand{\uu}{\mathfrak{U}}
\newcommand{\isou}{\mathcal{I}_\uu}
\newcommand{\dd}{\mathfrak{D}}
\newcommand{\isod}{\mathcal{I}_\dd}
\newcommand{\fp}{\overline{f(p)}}
\newcommand{\sss}{\mathfrak{S}}
\newcommand{\bb}[1]{\boldsymbol{#1}}
\newcommand{\bbl}[1]{\overline{\boldsymbol{#1}}}
\newcommand{\sn}{\hspace{0.08em} \lhd \lhd \hspace{0.3em}}
\DeclareMathOperator{\deft}{def}

\DeclareMathOperator{\Sl}{SL}
\newcommand{\frf}{\fr_\ff}
\newcommand{\isf}{\mathcal{I}_\ff}
\newcommand{\coho}{H^1(H,V)}
\newcommand{\iddots}{\reflectbox{$\ddots$}}

% THEOREMS ---------------------------------------------------------------
%%%%%%%%%\theoremstyle{plain}
\newtheorem{thm}{Theorem}
\newtheorem{cor}[thm]{Corollary}
\newtheorem{step}{Step} \newtheorem{lemma}[thm]{Lemma}
\newtheorem{prop}[thm]{Proposition} \newtheorem{rem}[thm]{Remark}
\newtheorem{num}[thm]{} \newtheorem{defn}[thm]{Definition}
\newtheorem{question}[thm]{Question} \newtheorem{con}[thm]{Conjecture}
\newtheorem{probl}[thm]{Problem}
\newtheorem{example}[thm]{Example}
\newtheorem{remark}[thm]{Remark}
\newtheorem{stepp}{Step}
%
\numberwithin{equation}{section}
\renewcommand{\theequation}{\thesection.\arabic{equation}}
\renewcommand{\footnote}{\endnote}
\newcommand{\ignore}[1]{}\makeglossary

%%%%%%%%%%%%%%%%%%%%%%%%%%%%%%%%%%%%%%%%%%%%%%%%%%%%%%%%%%%%%
%%%%%%%%%%%%%%%%%%%%%%%%%%%%%%%%%%%%%%%%%%%%%%%%%%%%%%%%%%%%%
%%%%%%%%%%%%%%%%%%%%%%%%%%%%%%%%%%%%%%%%%%%%%%%%%%%%%%%%%%%%%

\title[]{Profinite groups with many elements with large nilpotentizer and generalizations}

\author{Martino Garonzi}
\address{Martino Garonzi. University of Ferrara (Italy), Dipartimento di Matematica e Informatica.
ORCID: https://orcid.org/0000-0003-0041-3131}
\email{martino.garonzi@unife.it}

\author{Andrea Lucchini}
\address{Andrea Lucchini. University of Padova (Italy), Dipartimento di Matematica ``Tullio Levi Civita''. ORCID: https://orcid.org/0000-0002-2134-4991}
\email{lucchini@math.unipd.it}

\author{Nowras Otmen} 
\address{Nowras Otmen. University of Padova (Italy), Dipartimento di Matematica ``Tullio Levi Civita''. ORCID: https://orcid.org/0009-0009-8092-1689}
\email{nowras.naufel@math.unipd.it}

\thanks{Project funded by the EuropeanUnion - NextGenerationEU under the National Recovery and Resilience Plan (NRRP), Mission 4 Component 2 Investment 1.1- Call PRIN 2022 No. 104 of February 2, 2022 of Italian Ministry of University and Research; Project 2022PSTWLB (subject area: PE - Physical Sciences and Engineering) ``Group Theory and Applications''. The second author is member of GNSAGA (INDAM)}

\date{}

\subjclass[2020]{20E18}

\keywords{Profinite groups, pronilpotent groups, prosolvable groups}

\begin{document}
	\bibliographystyle{amsplain}

\begin{abstract}
Given a profinite group $G$ and a family $\mathcal{F}$ of finite groups closed under taking subgroups, direct products and quotients, denote by $\mathcal{F}(G)$ the set of elements $g \in G$ such that $\{x \in G\ |\ \langle g,x \rangle \ \mbox{is a pro-}\mathcal{F} \mbox{ group}\}$ has positive Haar measure. We investigate the properties of $\mathcal{F}(G)$ for various choices of $\mathcal{F}$ and its influence on the structure of $G$.
\end{abstract}

\maketitle

\section{Introduction}
Let $\mathcal F$ be a class of finite groups closed under taking subgroups, quotients and direct products. Given a profinite group $G$ and $x$ in $G$, let $\mathcal F_G(x)$ be the set of the elements $g\in G$ such that $\langle g, x\rangle$ (here and henceforth, 
for any subset $X$ of $G$, writing $\langle X\rangle$ we mean the smallest \emph{closed} subgroup of $G$ generated by $X$) is a pro-$\mathcal F$ group. Let $\mu_G$ be the normalized Haar measure of $G$. By \cite[Section 2]{aemp}, $\mathcal{F}_G(x)$ is closed in $G$, hence measurable, and 
$$\mu_G(\mathcal{F}_G(x)) = \inf_{N \unlhd_{\circ} G} \frac{|\mathcal{F}_{G/N}(xN)|}{|G/N|}.$$

For a positive real number $\varepsilon,$ let $\mathcal F_\varepsilon(G)$ be the set of the elements $g\in G$ such that $\mu_G(\mathcal F_G(g))\geq \varepsilon$. Given an open normal subgroup $N$ of $G$, let
$$\mathcal F_{\varepsilon,N}(G)=\{g\in G\mid \mu_{G/N}(\mathcal F_{G/N}(gN))\geq \varepsilon\}.$$ Notice that $\mathcal F_{\varepsilon,N}(G)$ is closed in $G$ and consequently $\mathcal F_\varepsilon(G)=\cap_N \mathcal F_{\varepsilon,N}(G)$ is closed. But then
$\mathcal F(G)=\cup_{n\in \mathbb N}\mathcal F_{1/n}(G)$ is a Borel subset of $G$ and coincides with the set of the elements $g\in G$ such that $\mu(\mathcal F_G(g))$ is positive.

The previous definitions give rise to the following questions:
\begin{question}\label{q1}Is there a group-theoretical characterization of the elements $g\in G$ with the property that $\mu(\mathcal F_G(g))$ is positive? 
\end{question}

\begin{question}\label{q2}Is $\mathcal F(G)$ a closed  subset of $G$ for every profinite group $G$?
\end{question}

\begin{question}\label{q3}Is $\mathcal F(G)$ a subgroup of $G$ for every profinite group $G$?
\end{question}

The answers to \cref{q2} and \cref{q3} are in general negative. We analyze \cref{q2} in further detail in \cref{families}. For \cref{q3}, if we take a class $\mathcal F$ that doesn't contain a cyclic group of order a power of a prime $p$, then this is already not true. Suppose $C_{p^m} \notin \mathcal F$ and consider $C_q \in \mathcal F$ with $q\neq p$. Let $G=P \rtimes C_q$ be a Frobenius group with $\exp(P)=p^m$. Then $\mathcal F(G)$ is not a subgroup of $G$, since if $x \in G$ has order $q$ then $x^P=xP$ so every element of $P$ is a product of two elements of order $q$.

In the particular case of the class $\mathcal A$ of the finite abelian groups, we can answer all three questions. Since $\mathcal A_G(g)=C_G(g)$ for every $g\in G$, the elements $g$ such that $\mu(\mathcal A_G(g))>0$ are those with finite conjugacy class. Then, it follows that $\mathcal A(G)$ coincides with the FC-center of $G$, which is a (not necessarily closed) subgroup of $G$. If $G=\prod_{i \in I} G_i$ is a direct product of groups, the corresponding restricted direct product is the subset of the direct product consisting of tuples $(g_i)_i$ such that $\{i \in I\ |\ g_i \neq 1\}$ is finite. An example of an FC-center that is not closed is the following: consider $G$ the direct product of infinitely many nonabelian simple groups, in which case the FC-center of $G$ coincides with the restricted direct product of the factors, which is dense but not closed in $G.$ In particular, if $\mu(\mathcal A(G))>0,$ then $\mathcal A(G)$ is an open normal subgroup of $G$ and it follows from \cite[Lemma 2.6]{sh} that the derived subgroup of $\mathcal A(G)$ is finite (which implies that  $\mathcal A(G)$ is virtually abelian). 

For most of the natural examples of classes $\mathcal F$, one can find at least one profinite group $G$ such that the set $\mathcal F(G)$ is not closed. Although we cannot fully characterize the families for which this occurs, the following result is a step in this direction.

\begin{thm}\label{allfingrps}
Let $\mathcal F$ be a class of finite groups closed under taking subgroups, quotients and direct products. Assume that $\mathcal F$ satisfies the following property: for every prime $p$, if $\mathcal F$ contains a finite nontrivial $p$-group, then it contains all the cyclic finite $p$-groups. Then $\mathcal{F}(G)$ is closed in $G$ for every profinite group $G$ if and only if $\mathcal F$ is the class of all finite groups.
\end{thm}

We also show in \cref{exp2} that, if the exponent of the groups in $\mathcal{F}$ is bounded, then we cannot hope to have the above result for $\mathcal{F}$.

\

Recall that the prosolvable radical of a profinite group $G$ is the subgroup of $G$ generated by all the closed normal prosolvable subgroups of $G$. It is denoted by $R(G)$ and it is a closed normal prosolvable subgroup of $G$. The following theorem provides a characterization of $\mathcal F(G)$ for $\mathcal F = \mathcal S$ the class of finite solvable groups.

%and provides, along with \cref{thodd}, a generalization of \cite[Theorem 1.2]{aemp}.

\begin{thm}\label{descs}
Let $\mathcal{S}$ be the class of finite solvable groups. If $G$ is a profinite group, then $\mathcal{S}(G)$ is a subgroup of $G$, $R(G) \leq \mathcal{S}(G)$ and $\mathcal{S}(G)/R(G)$ is the FC-center of $G/R(G)$. Moreover, if $\mathcal S(G)$ is closed, then $\mathcal S(G)/R(G)$ is finite.
\end{thm}





In \cite{aemp}, it is proven that if $\mathcal S(G) = G$, then $G$ is virtually prosolvable. We generalize this result in two different ways, both showing that you don't need \emph{every} element in $G$ to have a positive measure solvabilizer to guarantee virtual prosolvability.

\begin{thm}\label{musolpos} Let $G$ be a profinite group. If $\mu(\mathcal S(G)) > 0$, then $G$ is virtually prosolvable.
\end{thm}

\begin{thm}\label{thodd}
 Let $G$ be a profinite group. If $\mu(\mathcal S_G(g)) > 0$ for every $g$ of odd order (as supernatural number), then $G$ is virtually prosolvable. 
\end{thm}



Let $\mathcal{N}$ be the family of all finite nilpotent groups. Recall that the {(topological)} Fitting subgroup of a profinite group $G$ is the subgroup topologically generated by all the closed normal pronilpotent subgroups of $G$. It is denoted by $\fit(G)$ and it is a closed pronilpotent normal subgroup of $G$.

\begin{thm}\label{nilpp}
Let $G$ be a profinite group and let $g \in G$. If $\mu(\mathcal{N}_G(g)) > 0$, then there is an open subgroup $H$ of $G$ such that $g \in \fit(H)$. In particular, $g$ has finite order modulo $\fit(G)$.
\end{thm}


It follows from \cite[Theorem 1.3 (3)]{aemp} that, if $\mathcal N(G) = G$, then $G$ is virtually pronilpotent, thus its Fitting subgroup is open. The following result is a local version of this.

\begin{thm}\label{localnilp}
	Let $G$ be a  profinite group and let $P$ be a $p$-Sylow subgroup of $G.$  Suppose that $\mu(\nil_G(g))>0$ for every $g\in P.$
	Then the following hold:
	\begin{enumerate}
		\item $N_G(P)$ is an open subgroup of $G.$
		\item $P/O_p(G)$ is finite.
		\item $O_p(G)C_G(g)$ is open in $G$ for every $g \in P$.
		\item If $P$ is finite, then $C_G(P)$ is open in $G$.
	\end{enumerate}
\end{thm}

By \cite[Theorem 1.3]{aemp},  if $\mathcal N(G)=G$, then $G$ is virtually pronilpotent. In \cite{aemp}, it is observed that the condition $\mathcal N(G) = G$ is stronger than $G$ being virtually pronilpotent; indeed, it suffices to consider the example $\mathbb Z_p \rtimes C_2 = \mathbb Z_p \rtimes \la x \ra$, where $p$ is an odd prime and $x$ acts on the $p$-adic integers by inversion, and observe that $\mathcal N_G(x) = \{1,x\}$. It is natural to ask whether, as happens for the class $\mathcal S$ of finite solvable groups, we can reach the conclusion that $G$ is virtually pronilpotent from the  weaker assumption $\mu(\mathcal{N}(G)) > 0.$ The lack of a good characterization of the elements of $\mathcal N(G)$ in terms of the group structure, does not allow us to mimic the argument used in the proof of \cref{thodd} for the  class $\mathcal S$. Therefore, this is an apparently difficult question that we leave open. However, we obtain a partial result, providing an affirmative answer to the analogous question formulated for the class $\mathcal{P}_p(G)$ of finite $p$-groups, for any fixed prime $p.$

\begin{thm}\label{classp}
Let $G$ be a profinite group. If $\mu(\mathcal{P}_p(G)) > 0$ then $G$ is virtually pro-$p$.
\end{thm}


%It is interesting to ask whether $\mu(\mathcal{N}(G)) > 0$ implies that $G$ is virtually pronilpotent. This is a hard question because we do not have a group-theoretical characterization of $\mathcal{N}(G)$. This contrasts with the case of $\mathcal S(G)$; the characterization of elements in $\mathcal S(G)$ conjectured in \cite[Conjecture 1]{solv} and proved in \cite[Theorem 6]{solv} and \cite[Corollary 1.3]{fgg} is fundamental in the proof of \cref{descs}.

If $G$ is a finite group, then the intersection $\cap_{g\in G}\nil_G(g)$ of the nilpotentizers coincides with the hypercenter $Z_\infty(G)$ of $G$ (see \cite[Proposition 2.1]{az}). More generally, when $G$ is a profinite group, $\cap_{g\in G}\nil_G(g)$ is the hypercenter $Z_\infty(G)$, defined as the set of all $x \in G$ such that $xN \in Z_\infty(G/N)$ for every open normal subgroup $N$ of $G$ (see \cite{hyper}). In other words, if $g\in G,$
then $\nil_G(g)=G$ if and only if $g\in Z_\infty(G).$
So it is natural to ask whether \cref{nilpp} can be replaced with a stronger statement: {\sl{Suppose that $g\in G$ is such that $\mu(\nil_G(g))>0.$ Does there exist an open subgroup $H$ of $G$ such that $g\in Z_\infty(G)$?}}
%\textcolor{red}{The same argument used in the proof of \cref{nilpp} demonstrates that an affirmative answer to the previous question would imply that $g$ has finite order modulo $Z_\infty(G).$} \textcolor{blue}{Lo stesso argomento non funziona perché non è vero che, se $M\unlhd_o G$, allora $Z_\infty(M) \leq Z_\infty(G)$ (ad esempio $G=S_3$ e $M=A_3$).} The following result shows that this is not always the case.

\begin{thm}\label{nofiniteorder}
For every positive real number $\varepsilon$ there exists a solvable profinite group $G$ which contains an element $g$ such that {$\mu(\mathcal{N}_G(g)) > 1-\varepsilon$, $g$ has infinite order modulo $Z_{\infty}(G)$ and} there is no open subgroup $H$ of $G$ with $g \in Z_{\infty}(H)$.
\end{thm}

Let $G$ be a profinite group and let $x\in G.$ Consider the set $$\Lambda_G(x)=\{g\in G \mid \langle x, x^g\rangle \text { is pronilpotent}\}.$$ Clearly $\nil_G(x)\subseteq \Lambda_G(x)$, so it is reasonable to ask to what extent the results proven in  Section \ref{nilpot} remain true when substituting $\nil_G(x)$ with $\Lambda_G(x).$ By the Baer-Suzuki  Theorem (cf. \cite[p. 105]{gor}), if $\Lambda_G(x)=G,$ then $x \in \fit(G).$ A natural question is whether the following generalized version of \cref{nilpp} is true: {\sl{if $\mu(\Lambda_G(x))>0,$ then the order of $x$ modulo $\fit(G)$ is finite}.} 
We can state also a $p$-local version of the previous question: {\sl{Assume that $x$ is a $p$-element of $G$ and let $\Lambda_{G,p}(x)=\{g\in G \mid \langle x, x^g\rangle \text { is a pro-$p$ group}\}.$ Does $\mu(\Lambda_{G,p}(x))>0$ imply that the order of $x$ modulo $O_p(G)$ is finite?}}
We prove that the last question has in general a negative answer. This is implied by the following result.

\begin{thm}\label{nobaer}
For every positive real number $\eta<1$, every integer $t \geq 1$
and every prime $p,$ there exists a finite group $G$ such that $O_p(G)=1$ and
$|\Lambda_{G,p}(x)|\geq \eta|G|$ for some element $x \in G$ of order $p^t$.
\end{thm}

Indeed, for every positive integer $t$, choose a finite group $G_t$ with $O_p(G_t)=1$ and containing an element $x_t \in G_t$ of order $p^t$ such that $|\Lambda_{G_t,p}(x_t)| \geq \eta_t |G_t|$ where $\eta_t = 1-1/2^t$. Let $G =\prod_{t \geq 1} G_t$ and $x=(x_t)_t \in G$. It follows that
\[\mu(\Lambda_{G,p}(x)) \geq \prod_{t \geq 1} \eta_t > 0.\]
Moreover $O_p(G)=1$ and $x$ has infinite order.

In \cref{families} we prove \cref{allfingrps}. In \cref{solv} we prove Theorem \ref{descs}, \ref{musolpos} and \ref{thodd}. In \cref{nilpot} we prove Theorems  \ref{nilpp}, \ref{localnilp}, \ref{classp} and \ref{nofiniteorder}. In \cref{Lambda} we prove \cref{nobaer}.

\section{General families of finite groups} \label{families}

In this section we discuss \cref{q2}. In particular we prove Theorem \ref{allfingrps} and, in \cref{exp2}, we give an example to show that the statement of this theorem does not hold if we remove the assumption.

\begin{proof}[Proof of \cref{allfingrps}]

If $\mathcal F$ is the class of all finite groups, then $\mathcal{F}(G)=G$ for every profinite group $G$ and this proves the implication $(\Leftarrow)$.

Now we prove the other implication. Assume that $\mathcal F$ is not the class of all finite groups. We will construct a profinite group $G$ such that $\mathcal{F}(G)$ is not closed.

Assume first that $\mathcal F$ contains all the finite cyclic groups. Since $\mathcal F$ is not the class of all finite groups and every finite group is isomorphic to a subgroup of an alternating group, there exists an alternating group $\alt(m)$ which does not belong to $\mathcal F$. Since $\mathcal{F}$ contains all finite cyclic groups, $\mu(\mathcal{F}_{\alt(m)}(1))=1$. It is known that, for every $x \in \alt(m) \setminus \{1\}$, there exists $y \in \alt(m)$ such that $\langle x,y \rangle = \alt(m)$ \cite{gk}, therefore $0< \mu(\mathcal{F}_{\alt(m)}(x)) < 1$ for all $x \in \alt(m) \setminus \{1\}$. Let $G := \alt(m)^{\mathbb{N}}$ and let $c := \max_{x \in \alt(m)} \mu(\mathcal{F}_{\alt(m)}(x))$. Then $0 < c < 1$. If $g = (g_n)_n \in G$ then
$$\mu(\mathcal{F}_G(g)) = \prod_{n \in \mathbb{N}} \mu(\mathcal{F}_{\alt(m)}(g_n)) = \prod_{g_n \neq 1} \mu(\mathcal{F}_{\alt(m)}(g_n)) \leq \prod_{g_n \neq 1} c.$$
It follows that $\mu(\mathcal{F}_G(g)) > 0$ if and only if $g_n$ is almost always $1$, so $\mathcal{F}(G)$ equals the restricted direct power $\alt(m)^{(\mathbb{N})}$, and hence it is not closed in $G$ (it is dense).

Now assume that $\mathcal F$ does not contain all the finite cyclic groups. By our assumption on $\mathcal F$, there exist two prime numbers $p,q$ such that $C_p$ is in $\mathcal F$ and $C_q$ is not in $\mathcal F$. Let $n$ be the order of $q$ in the multiplicative group of $\mathbb{Z}/p\mathbb{Z}$, so that $p$ divides $q^n-1$. Let $Q$ be the additive group of the field with $q^n$ elements. The group $P \cong C_p$ generated by an element of multiplicative order $p$ in the field with $q^n$ elements acts on $Q$ fixed point freely (by multiplication), so $H=Q \rtimes P$ is a Frobenius group. Fix a positive integer $t$ and let $n_t=p^t$, $\sigma=(1,\ldots,n_t) \in S_{n_t}$, $h \in H \setminus Q$, $Y_t= Q^{n_t} \langle g \rangle < H \wr \langle \sigma \rangle$ where $g=(h,1,\ldots,1)\sigma$. Note that $|Y_t| = q^{p^t n} p^{t+1}$ and that a subgroup of $Y_t$ belongs to $\mathcal{F}$ if and only if it is a $p$-subgroup.
We want to prove that, if $u \in Y_t$, then $\mu(\mathcal F_{Y_t}(u)) \leq q^{-{n_tn}}.$ %being $\mu(\mathcal{F}_{Y_t}(u))$
%the probability that a randomly chosen $w \in Y_t$ is such that $\langle u, w\rangle \in \mathcal{F}$.
Notice that $\langle g\rangle $ is a $p$-Sylow subgroup of $Y_t$ and that $g^{n_t}=(h,\dots,h)$ generates the unique subgroup of order $p$ of $\langle g \rangle.$ Moreover
$(y_1,\dots,y_{n_t})\in Q^{n_t}$ normalizes $\langle g^{n_t}\rangle $ if and only if $y_i=1$ for every $1\leq i\leq n_t$ and therefore $\langle g\rangle$ is the unique $p$-Sylow subgroup of $Y_t$ containing $g^{n_t}.$ This implies that  $\langle g^{n_t}, x\rangle \in \mathcal{F}$ if and only if $x \in \langle g\rangle$ and therefore $$\mu(\mathcal{F}_{Y_t}(g^{n_t}))=\frac{|g|}{|Y_t|}=\frac{1}{q^{n_tn}}.$$
Let $u$ be an arbitrary nontrivial $p$-element of $Y_t$. 
Since $\mu(\mathcal{F}_{Y_t}(u))\leq \mu(\mathcal{F}_{Y_t}(u^m))$ for every $m\in \mathbb N,$ we may assume $|u|=p$. But then $u$ is conjugate to an element of $\langle g^{n_t} \rangle$, and therefore
$$\mu(\mathcal F_{Y_t}(u))\leq \frac{|g|}{|Y_t|}=\frac{1}{q^{n_tn}}.$$
Now consider $G=\prod_{t\in \mathbb N}Y_t$.
Assume $y=(y_t)_{t\in \mathbb N}\in G$ and let $\Lambda:=\{t\mid y_t\neq 1\}.$ Then
$$\mu(\mathcal F_{G}(y))\leq \prod_{t\in \Lambda}\frac{1}{q^{n_tn}},$$ so if $\mu(\mathcal{F}_{G}(y))>0$, then $|\Lambda|$ is finite (and clearly $y_t$ is a $p$-element for every $t\in \Lambda$). Conversely, suppose that there exists a finite $\Lambda \subset \mathbb{N}$ such that $y_t=1$ if $t\notin \Lambda,$ $y_t$ is a $p$-element of $Y_t$ otherwise. If $t\in \Lambda,$ then $\alpha_t:=\mu(\mathcal{F}_{Y_t}(y_t))>0$. On the other hand, since every element of $Y_t \setminus Q^{n_t}$ is a $p$-element, $$\mu(\mathcal{F}_{Y_t}(1)) \geq \frac{|Y_t|-|Q^{n_t}|}{|Y_t|} = 1-\frac{1}{p^{t+1}}.$$ Hence
$$\mu(\mathcal F_{G}(y))\geq \prod_{t\in \Lambda}\alpha_t
\prod_{t\notin \Lambda}\left(1-\frac{1}{p^{t+1}}\right)>0.$$
We have proved that  $\mathcal{F}(G)$ equals the {set of $p$-elements belonging to the restricted direct product of the groups $Y_t$}. Any $p$-element in $G$ belongs to the closure of $\mathcal{F}(G)$, so it is not closed. 
\end{proof}

On the other hand, we are able to find a family $\mathcal F$ for which $\mathcal F(G)$ is closed for \emph{every} profinite group $G$. {If $p$ is a prime number, let $\Omega_p(G)$ be the set of $p$-elements of $G$.}

\begin{prop} \label{exp2}
Let $\mathcal F$ be the class of all finite groups of exponent at most $2$. Then $\mathcal{F}(G)$ is closed in $G$ for every profinite group $G$.
\end{prop}
\begin{proof}
If $x,y \in G$ are such that $y \in \mathcal{F}_G(x)$ then $x,y$ have order at most $2$ and $xy=yx$. In particular, if $x \in \mathcal{F}(G)$, i.e. $\mu(\mathcal{F}_G(x))>0$, then the probability that an element of $G$ has order $2$ is positive. So we may assume that $\mu(\Omega_2(G))>0$, since otherwise $\mathcal{F}(G) = \varnothing$ is closed. By \cite[Theorem 2]{LP}, $G$ is abelian-by-finite, i.e. there exists an abelian open normal subgroup $N$ of $G$. Write $G=t_1N \cup \ldots \cup t_rN$ with $r=|G:N|$. If $t_iN$ contains elements of order $2$, then we may assume $t_i$ has order $2$ and the set of elements of order $2$ in $t_iN$ coincides with $t_iH_i$ where $H_i=\{n \in N\ |\ n^{t_i}=n^{-1}\} \leq N$. We assume that $t_i$ has order $2$ precisely when $1 \leq i \leq s$, where $s \leq r$. We want to prove that $\mathcal{F}(G)$ is closed. Since $\mathcal{F}(G)$ consists of elements of order at most $2$, it is enough to prove that $\mathcal{F}(G) \cap t_jH_j$ is closed for all $j=1,\ldots,s$. Without loss of generality, it is enough to prove that $Y=\mathcal{F}(G) \cap t_1H_1$ is closed. Let $x=t_1h_1 \in Y$ where $h_1 \in H_1$. There must exist $j$ such that $\mu(\mathcal{F}_G(x) \cap t_jH_j) > 0$. Let $$\Lambda_j := \{x \in Y\ |\ \mu(\mathcal{F}_G(x) \cap t_jH_j) > 0\}.$$
Then $Y=\Lambda_1 \cup \ldots \cup \Lambda_s$, so it is enough to prove that $\Lambda_j$ is closed for each $j=1,\ldots,s$. Assume $x = t_1h_1 \in \Lambda_j$. Note that the fact that $\Lambda_j \neq \varnothing$ implies that $H_j$ is a closed subgroup of $G$ with positive Haar measure, so $|G:H_j|$ is finite. If $t_jh_j \in t_jH_j$, then saying $t_jh_j \in \mathcal{F}_G(x)$ is equivalent to writing $t_1h_1t_jh_j=t_jh_jt_1h_1$. Since $H_1$ and $H_j$ are contained in $N$, which is abelian, we can restate this condition as $t_1 t_j^{h_1^{-1}} = t_j t_1^{h_j^{-1}}$. Moreover, $H_j \cap C_G(t_1h_1) = C_{H_j}(t_1)$ so $t_j^{-1} (\mathcal{F}_G(x) \cap t_jH_j)$ is a coset of $C_{H_j}(t_1)$ in $H_j$. So if $|G:C_{H_j}(t_1)|$ were infinite, then $\mu(\mathcal{F}_G(x) \cap t_jH_j)=0$ would imply $x \not \in \Lambda_j$, a contradiction. Therefore, if $\Lambda_j \neq \varnothing$, then $|G:C_{H_j}(t_1)|$ is finite. Thus, on one hand, we can conclude that all elements $t_1h_1 \in \Lambda_j$ are contained in $X_j =\{t_1h_1 : h_1 \in H_1 \text{ and }t_1t_j^{h_1^{-1}} \in t_jt_1^{H_j}\}$. On the other, if $x = t_1h_1 \in X_j$, from $\mu(\mathcal{F}_G(x) \cap t_jH_j) = \mu(C_{H_j}(t_1)) = 1/|G:C_{H_j}(t_1)| > 0$ it follows that $x$ is in $\Lambda_j$. Therefore, $\Lambda_j$ equals one of $\varnothing$, $X_j$, hence it is closed.
\end{proof}

Notice that the above proof relies heavily on the fact that, if $G$ satisfies the identity $x^2=1$ with positive probability, then it is virtually abelian. It is natural to wonder what happens for a class $\mathcal F$ of finite groups with uniformly bounded exponent, but we can't use the same strategy as above since we lack a similar structural result for profinite groups which satisfy $x^n=1$ with positive probability for $n>2$.

\begin{question}
    Let $\mathcal F$ be a class of finite groups closed under taking subgroups, quotients and direct products and assume that it has uniformly bounded exponent. Is there a profinite group $G$ such that $\mathcal F(G)$ is not closed?
\end{question}

\section{The class of finite solvable groups}\label{solv}

In this section, we discuss the case in which $\mathcal F$
coincides with the class $\mathcal S$ of the finite solvable groups. The set $\mathcal S_G(x)$ consists of all the  elements $y\in G$ such that the subgroup generated by $x$ and $y$ is prosolvable and is called the \emph{solvabilizer} of $x$ in $G.$ The characterization of elements with solvabilizer  of positive measure was proposed in \cite[Conjecture 1]{solv} and proved in \cite[Theorem 6]{solv} and \cite[Corollary 1.3]{fgg}. {The authors of these papers} obtained the following: given a profinite group $G$ and a positive integer $n$, let $\Sigma_n(G)$ be the set of the elements $g$ in $G$  that centralize all but at most $n$ nonabelian factors in a chief series of $G/N$ for every open normal subgroup $N$ of $G.$ Then, we have that $\mathcal S(G)=\cup_{n\in \mathbb N}(\Sigma_n(G))$, answering \cref{q1} and \cref{q3}: $\mathcal{S}(G)$ is a subgroup of $G$ which admits a group-theoretical characterization. For \cref{q2}, it suffices to use \cref{allfingrps} to deduce that $\mathcal{S}(G)$ is not closed in $G$ in general.

Our first aim in this section is to use the aforementioned characterization and obtain a better understanding of the structure of $\mathcal S(G).$ If $G$ is a profinite group, then we will denote by $R(G)$ the largest subgroup of $G$ which is closed, normal and prosolvable. It is well known that $g\in R(G)$ if and only if $gN\in R(G/N)$ for every open normal subgroup $N$ of $G.$

\begin{lemma}
Let $G$ be a finite group. If $g$ centralizes all the nonabelian chief factors of $G$, then $g\in R(G).$
\end{lemma}
\begin{proof}
	We prove the statement by induction on $|G|.$ Let $N$ be a minimal normal subgroup of $G$ and let $K/N=R(G/N)$. By induction $g\in K.$ If $N$ is abelian, then $K=R(G)$ and we are done. If $N$ is nonabelian, then $C_K(N)\cap N=Z(N)=1.$ Thus $C_K(N)$ is solvable, and therefore $C_K(N)\leq R(K)\leq R(G).$ Since $g\in C_K(N),$ we are done.
\end{proof}

From now on, if $G$ is a profinite group and $g \in G$, we will write that $g$ centralizes all the nonabelian chief factors of $G$ to mean that $gM$ centralizes all the nonabelian chief factors of $G/M$ for every open normal subgroup $M$ of $G.$

\begin{cor}\label{corradical}
Let $G$ be a profinite group. If $g \in G$ centralizes all the nonabelian chief factors of $G$, then $g\in R(G).$
\end{cor}



\begin{lemma}\label{fccenter}Let $G$ be a profinite group. Then $R(G) \leq \mathcal S(G)$. Moreover $\mathcal S(G)/R(G)$ is contained in the FC-center of $G/R(G)$ and all elements of $\mathcal S(G)/R(G)$ have finite order.
\end{lemma}
\begin{proof}
For every $g \in R(G)$ and every $x \in G$, we have that $\la x,g\ra \leq R(G)\langle x \rangle$, which is prosolvable, so $\mathcal S_G(g) = G$ and $R(G) \leq \mathcal S(G)$. Thus is not restrictive to assume $R(G)=1$. Let $g\in \mathcal S(G).$ There exists an open normal subgroup $N$ of $G$ such that $g$ centralizes all the nonabelian chief factors of $N.$ Let $K=N\langle g\rangle$. Then, by \cref{corradical}, $g\in R(K).$ Since $R(G)=1$, we must have $R(N)=1$. In particular this implies $R(K)\cap N=1$, so $[R(K),N] = 1$. It follows that $N\leq C_G(g)$, therefore $g$ is contained in the $FC$-center of $G$ . Moreover $g$ has finite order since $\langle g\rangle \cap N = 1$.
\end{proof}

\begin{lemma}\label{inclus}Let $G$ be a profinite group. Then the $FC$-center of $G$ is contained in $\mathcal S(G).$
\end{lemma}

\begin{proof}
Let $g$ be an element of the $FC$-center of $G$. Then $C_G(g)$ is open and therefore it contains an open normal subgroup $N$ of $G$. Clearly $g$ centralizes all the nonabelian chief factors of $G$ contained in $N$ and therefore $g\in \mathcal S(G).$
\end{proof}

Recall that if $G$ is a profinite group, $N \unlhd_c G$ and $\pi \colon G \to G/N$ is the natural projection, then $\mu_{G/N}(S) = \mu_G (\pi^{-1}(S))$ for every measurable subset $S$ of $G/N$. If $R=R(G)$, then one can prove that $\pi^{-1}(S_{G/N}(gR)) = S_G(g)$ and so $\mathcal S(G/R(G)) = \mathcal S(G)/R(G)$.

\begin{proof}[Proof of \cref{descs}]
The first part of the statement  follows from Lemmas \ref{fccenter} and \ref{inclus}. Furthermore, if $\mathcal S(G)$ is closed, then $\mathcal S(G)/R(G)$ is a profinite FC-group and contains an abelian open normal subgroup $N/R(G)$ by \cite[Lemma 2.6]{sh}. The preimage of $N$ is open and prosolvable in $\mathcal S(G)$, so $\mathcal S(G)/R(G)$ is finite.
\end{proof}

Note that when $R(G)=1$ every element in $\mathcal S(G)$ has finite order. This immediately gives examples of profinite groups $G$ such that $\mathcal S(G) = 1$; it suffices to consider torsion-free profinite groups with trivial prosolvable radical (\emph{e.g.}, non-abelian free profinite groups). For examples of groups $G$  with $\mathcal S(G)=1$ but containing non-trivial elements of finite order, one may look at the hereditarily just-infinite, non-virtually prosolvable groups obtained by infinitely iterated wreath products (see \cite{mv}).

%The following corollary provides examples of profinite groups $G$ with trivial $\mathcal S(G)$.


%\begin{cor}If $F$ is a nonabelian free profinite group, then $\mathcal S(F)=1.$\end{cor}
%\begin{proof}We may deduce from \cite[Theorem 1]{pop} that if $g$ is a free generator of $F,$ then $\mathcal S_F(g)={\langle g\rangle},$ so$R(F)=\cap_{x\in F}\mathcal S_F(x)=1$, since the intersection of the procyclic   subgroups generated by the free generators is in the center of $F$, which is trivial. Hence, by \cref{descs}, $\mathcal S(F)$ coincides with the FC-center of $F.$ However, if $g$ is an element of the FC-center of $F$, then $C_F(g)$, being open, is a nonabelian free profinite group. Hence $g\in Z(C_F(g))=1.$\end{proof}

%Using \cref{descs}, one can prove that many other profinite groups $G$ have a trivial $\mathcal S(G)$, for instance, just-infinite non-virtually prosolvable profinite groups and non-virtually prosolvable groups that split as a non-trivial free profinite product of profinite groups.
% {\color{blue}{Poco convinto sull'opportunità del corollario 19. Intanto è falso per $n=1$. Poi penso che $F$ non possa essere virtualmente provolvable, eccetto il caso $C_2\coprod C_2$ che è il completamento profinito del diedrale infinito. Non sono convinto che serva la teoria di Bass-Serre per dire che $R(G)=1$. Non viene facile dal teorema di Pop? Ho paura che tiriamo in ballo cose complicate per risultati piccoli, non cosi strettamente pertinenti al lavoro}} {\color{red} Certo, si intende che $n$ sia più grande di 1. Sì, il $C_2 \amalg C_2$ è l'unico virtualmente prorisolubile, però l'ho scritto così perché mi sembrava più semplice. Comunque, in generale non penso che il lavoro di Pop sia abbastanza, dato che non dice niente sui gruppi normali prorisolubili. Se preferite toglierlo, togliamolo, sono d'accordo che le cose qui sono molto più diverse da quelle che usiamo nel resto del lavoro.}
% {\color{red} I think I can improve the above corollary. I might be able to improve it further for a free profinite product over any profinite space, but I'm not fully sure how to do it still.
% \begin{cor}
%     Let $G = \coprod_{i=1}^{n} G_i$ be the free profinite product of profinite groups $G_i$ and suppose that $G$ is not virtually prosolvable. Then $\mathcal S(G)=1$.
% \end{cor}
% \begin{proof}
%     We use results of profinite Bass-Serre theory. Note that $G$ is the fundamental profinite group of a reduced injective graph of profinite groups over a finite graph (\cite[Example 6.2.3 (b)]{rib}) and, in particular, since it's not virtually prosolvable, it's also not of \emph{special type}. Then, we are under the hypothesis of \cite[Theorem 7.2.6]{rib}, which says that a normal subgroup $N\unlhd_c G$ is either contained in every edge group, which are trivial in our case, or $N$ contains an open subgroup $H$ which surjects onto a free nonabelian profinite group. From this, it's immediate that $R(G) = 1$. Next, take an element $g$ in the FC-center of $G$. Since $C_G(g)$ is open, then $C_G(g)$ splits as a free profinite product by the Kurosh theorem for open subgroups of free profinite products of finitely many profinite groups (\cite[Theorem 7.3.1]{rib}). Note that $\la g \ra \unlhd_c C_G(g)$, so, using \cite[Theorem 7.2.6]{rib} once more, there cannot be an $H \leq_o \la g \ra$ that surjects onto a nonabelian free profinite group, so $\la g \ra$ must be contained in every edge group, all of which are trivial. So $g = 1$ and by \cref{descs} it follows that $\mathcal S(G) = 1$.
% \end{proof}
% This applies, in particular, to finitely generated nonabelian free profinite groups, though a similar strategy yields the result for arbitrary free profinite groups.}

% Recall that a profinite group $G$ is called \emph{just-infinite} if every non-trivial closed normal subgroup of $G$ has finite index in $G$.

% \begin{cor}
%     If $G$ is just-infinite and non-virtually prosolvable, then $\mathcal S(G)=1$.
% \end{cor}
% \textcolor{blue}{(Martino) Scusate la domanda stupida, ma esistono gruppi profiniti just infinite e non virtualmente prorisolubili? (Andrea) Si, prodotti intrecciati reiteratici di gruppi alterni si comportano cosi. (Nowras) Secondo voi, va aggiunto dopo la dimostrazione che questi sono esempi di gruppi che soddisfano le ipotesi del corollario?} \textcolor{red}{(Martino) Secondo me s\`i}
% \begin{proof}
%     The assumptions easily imply that $R(G)=1$. It remains to show that the FC-center of $G$ is trivial. Suppose $x$ is an element in the FC-center of $G$. If $C_x$ is the conjugacy class of $x$, then $\la C_x \ra \lhd_{\circ} G$ and $Z =Z(\la C_x \ra) = (C_G(x))_G \cap \la C_x\ra\neq 1$, which holds because $C_G(x)$ has finite index and thus its core is non-trivial and open. Since $Z$ is characteristic in $\la C_x \ra$, it follows that $Z \lhd_{\circ} G$. Since $Z \leq R(G),$
%     it follows that $G$ is finite, a contradiction.
% \end{proof}

\

 Before proving \cref{musolpos}, we need a preliminary lemma.

% \begin{thm}\label{sigmapos}
% 	Let $G$ be a profinite group. If $\mu(\mathcal S(G))>0,$ then $G$ is virtually prosolvable, and consequently $\mathcal S(G)=G.$
% \end{thm}
% Before proving this theorem, we need a preliminary lemma.

\begin{lemma}\label{subseries}
	Let $G$ be a finite group, $H$ a normal subgroup of $G$ and $N$ a nonabelian minimal normal subgroup of $G$ contained in $H$. If $M$ is a normal subgroup of $H,$ then an $H$-chief series of $NM/M$ has length at most $|G:H|.$
\end{lemma}

\begin{proof}
	We have that $N=S_1\times \dots \times S_u,$ where $S_1,S_2,\dots,S_u$ are isomorphic nonabelian simple groups.
	Since $N$ is a minimal normal subgroup of $G$, $G$ permutes transitively the factors $\{S_1,\dots,S_u\}.$ Since $H$ is normal in $G,$ the orbits of $H$ on $\{S_1,\dots,S_u\}$ have all the same size, and the number of these orbits is at most $t=|G:H|.$ Notice that $NM/M \cong_H N/N\cap M$ and a minimal $H$-subgroup of $N/N\cap M$ corresponds to an $H$-orbit on $\{S_1,\dots,S_u\}.$ Thus an $H$-chief series of $NM/M$ has length at most $t.$
\end{proof}
 
\begin{proof}[Proof of \cref{musolpos}]
	Assume $\mu(\mathcal S(G))>0.$ Since $\mathcal S(G)$ is a subgroup of $G$, we must have that $\mathcal S(G)$ has finite index in $G$. Moreover by \cite[Lemma 1.1 (2)]{bm} $\mathcal S(G)$ is closed in $G.$ Let $t=|G:\mathcal S(G)|$  and let $g\in \mathcal S(G)$. Then there exists $m\in \mathbb N$, such that, for every open normal subgroup $N$ of $G,$ $g$ centralizes all but $m$ nonabelian factors in any chief series of $G/N.$ Now let $M$ be an open normal subgroup of $\mathcal S(G)$. The normal core $K$ of $M$ in $G$ is open in $G$. We may so find a chief series
	$1=N_0/K < \dots < N_u/K=\mathcal S(G)/K<\dots  < N_v/K=G/K$ of $G/K$. Then
	$N_0M/M \leq \dots \leq N_uM/M=\mathcal S(G)/M$ is a normal chain of $\mathcal S(G)/M.$ Notice that $N_iM/N_{i+1}M$ is not necessarily a chief factor of $\mathcal S(G)/M$, but by Lemma \ref{subseries} if $N_i/N_{i+1}$ is nonabelian, then the length of a $\mathcal S(G)$-chief series of
	$NM/M$ is at most $t$. Moreover if $g$ centralizes $N_iM/N_{i+1}M$, then it centralizes all the factors of a $\mathcal S(G)$-chief series of $N_iM/N_{i+1}M.$ Since $g$ centralizes all but $m$ of the nonabelian factors $N_i/N_{i+1}$, we conclude that $g$ centralizes all but $m\cdot t$  nonabelian factors of a chief series of $\mathcal S(G)/M.$ We have so proved that $g\in \mathcal S(\mathcal S(G)),$ that is, $\mathcal S(\mathcal S(G))=\mathcal S(G).$ By \cite[Theorem 1.2]{aemp}, $\mathcal S(G)$ is virtually prosolvable. But then, since $\mathcal S(G)$ is open in $G,$ $G$ also is virtually prosolvable.
\end{proof}

%Another generalization of \cite[Theorem 1.2]{aemp} is the following, done by considering solely the solvabilizer of elements of odd order.


We need a preliminary lemma for the proof of \cref{thodd}.

\begin{lemma}\label{oddele}
Let $N$ be a nonabelian minimal normal subgroup of a finite group $G$. If $g\in G$ and the order of $g$ modulo $N$ is odd, then there exists $n\in N$ such that $|gn|$ is odd and $gn \notin C_G(N).$
\end{lemma}

\begin{proof}
We may write $g=g_1g_2$ with $|g_1|$ odd, $|g_2|$ a 2-power and $[g_1,g_2]=1.$ Since $g$ has odd order modulo $N,$ $g_2\in N.$ So $g_1\in gN.$ This means that we may assume that $|g|$ is odd. If $g\notin C_G(N)$ we are done. Assume that $g\in C_G(N)$. In this case let $n$ be a non-trivial element of odd order in $N$. Then $|gn|$ is odd, and $gn\notin C_G(N)$ (since $n\notin C_G(N))$.
\end{proof}

%\begin{proof}[Proof of \cref{thodd}]
%With iterated applications of Lemma \ref{oddele}, arguing as in the proof of \cite[Corollary 22]{pel}, we may construct an element $g\in G$ of odd order and which centralizes no nonabelian chief factor of $G.$ By assumption, the solvabilizer of $g$ has positive Haar measure and therefore $g$ centralizes almost all the nonabelian chief factors of $G$. It follows that a chief series of $G$ \textcolor{blue}{(Martino) Abbiamo definito cos'\`e una chief series di un gruppo profinito?} \textcolor{green}{(Nowras) No, e penso che non sia possibile definirlo bene per gruppi profiniti qualsiasi. La riscrivo con più parole per essere più chiaro.} contains only finitely many nonabelian factors and therefore $G$ is virtually prosolvable.
%\end{proof}

\begin{proof}[Proof of \cref{thodd}]
We prove the contrapositive. Arguing as in the proof of \cite[Corollary 22]{pel} and using \cref{oddele}, we construct an element $g\in G$ of odd order and which centralizes no nonabelian chief factor of $G$. We repeat the argument as follows. If $G$ were not virtually prosolvable, we would be able to find a descending chain of open normal subgroups
\[G \geq N_1 > M_1 \geq N_2 > M_2 \geq \ldots\]
such that $N_i/M_i$ is a nonabelian minimal normal subgroup of $G/M_i$. It is not restrictive to suppose that $\cap_{i\in \mathbb N} M_i = 1$, since for any $g \in G$ and $K\unlhd_c G$, we have $\mu(\mathcal S_G(g)) \leq \mu(\mathcal S_{G/K}(gK))$. Thus, we may assume $G \cong \varprojlim_{i \in \mathbb N} G/M_i$ and the elements of $G$ will be of the form $(g_iM_i)_{i \in \mathbb N}$ such that $g_iM_i = g_jM_i$ whenever $j\geq i$. Take an element $g_1 \in G$ of odd order and consider $g_1M_2$. Its order modulo $N_2/M_2$ is odd, so by \cref{oddele} we may can choose $g_2 \in g_1N_2$ such that $g_2M_2$ does not centralize $N_2/M_2$. It is clear that $g_2 M_1 = g_1M_1$. By induction, we obtain elements $g_iM_i$ such that $g=(g_iM_i)$ is an element of $G$. Further, for every positive integer $n$, we can find $M\unlhd_o G$ such that 
$g$ does not centralize at least $n$ nonabelian chief factors of $G/M$. It follows that $\mu(\mathcal S_G(g))=0$.
\end{proof}


% Finally, instead of considering the class $\mathcal S$ of the finite solvable groups, we may consider the subclass $\mathcal S_{\rm{odd}}$ consisting of the finite groups of odd order. If $G$ is a finite group, then 
% $\mathcal S_{\rm{odd}}(G)$
% consists of the elements of odd order, which is not necessarily a subgroup of $G$, so in general the answer to \cref{q3} is negative. Notice that the only known examples such that $\mathcal F(G)$ is a subgroup are for the classes of finite solvable groups and of finite abelian groups. One could have then hoped that, for extension closed varieties, it would hold that the answer to \cref{q3} would be affirmative; however, $\mathcal S_{\rm{odd}}$ provides a counterexample for that.

%Let $G=\prod_{n \geq 2}{\rm{SL}}(2,2^n).$ \textcolor{blue}{$\leftarrow$ Questo gruppo \`e un controesempio per cosa?}\textcolor{red}{(Andrea) qui va tutto cambiato, questo esempio era nato come esempio in cui $F(G)$ non è chiuso, ma adesso non serve piu, va solo messo come controesempio a Larsen-Shalev, quindi va accorciato e scritto dopo di aver raccontato del rsultato di Larsen Shalev}
%Since $\mathcal S_{\rm{odd}}(G) \subseteq \mathcal S$
%and $R(G)=1,$ it follows from \cref{descs} that 
%$\mathcal S_{\rm{odd}}(G) \subseteq \mathcal S$ is contained in the $FC$-center of $G$, which coincides with the restricted direct product of the factors.
%Moreover if $g=(g_n)_{n\in \mathbb N}\in \mathcal S_{\rm{odd}}(G)$, then $|g_n|$ is odd for every $n\in \mathbb N$ (otherwise $\mathcal S_{{\rm{odd}},_G}(g) =\varnothing).$ Thus 
%$S_{\rm{odd}}(G)$ is contained in the subset
%$\Omega_{\rm{odd}}(G)$ consisting of the elements $(g_n)_{n\in \mathbb N}$
%such that $|g_n|$ is odd for every $n\in \mathbb N$
%and $g_n=1$ for all but finitely many $n\in \mathbb N.$ We claim that the equality $S_{\rm{odd}}(G)=\Omega_{\rm{odd}}(G)$ holds. Notice that the order of an arbitrary element of  ${\rm{SL}}(2,2^n)$ is either 2 or odd. Moreover the number of 2-elements in ${\rm{SL}}(2,2^n)$ is precisely $(2^n)^2$ (see for example \cite{moller}). One of the elements is the identity, so ${\rm{SL}}(2,2^n)$ has $(2^n-1)2^n(2^n+1)$ elements and $(2^n-1)2^n(2^n+1)-(2^n-1)(2^n+1)$ of these have odd order. So the probability that an element of ${\rm{SL}}(2,2^n)$ has odd order is precisely $1-\frac{1}{2^n}.$
%Then the probability that a random chosen element of $G$ has odd order (as a supernatural number) is
%$$\eta:=\prod_{n \in \mathbb N}\left(1-\frac{1}{2^n}\right)>0.$$
%Now let $g=(g_n)_{n\in \mathbb N}\in \Omega_{\rm{odd}}(G)$ and let $\Lambda(g)=\{n\in \mathbb N\mid g_n\neq 1\}.$
%For every $n \in \Lambda(g),$ let $p_n$ be the probability that an arbitrary element $x\in {\rm{SL}}(2,2^n)$ satisfies the property that $\langle x, g_n\rangle$ has odd order. Then
%$$\mu(\mathcal S_{{\rm{odd}},_G}(g))=\eta \prod_{n\in \Lambda(g)}\frac{p_n}{1-\frac{1}{2^n}}>0.$$
%We have so proved that $S_{\rm{odd}}(G)=\Omega_{\rm{odd}}(G)$ 
%and therefore $S_{\rm{odd}}$ is not closed in $G$ (indeed
%$\overline{S_{\rm{odd}}(G)}=\{(g_n)_{n\in \mathbb N}\mid |g_n| \text { odd for every $n\in \mathbb N$}\}).$

% Larsen and Shalev \cite{ls} proved the following result: {\sl{
% Let $G$ be a profinite group and suppose that the set of elements of $G$ of finite
% odd order has positive Haar measure. Then $G$ has a prosolvable open normal
% subgroup.}} This motives the following question:	{\sl{Let $G$ be a profinite group and suppose that the set of elements whose order, as a supernatural number, is odd has positive Haar measure. Is $G$ virtually prosolvable?}}
% The previous example provides a negative answer. Indeed $G=\prod_{n \geq 2}{\rm{SL}}(2,2^n)$ is not virtually prosolvable while the probability that a randomly chosen element of $G$ has odd order (as a supernatural number) coincides with the positive real number $\eta.$

%{\color{blue} (Nowras) I propose the following rewriting of the end of this section.

Finally, we make a remark regarding a result of Larsen and Shalev. In \cite{ls}, they prove the following: {\sl{Let $G$ be a profinite group and suppose that the set of elements of $G$ of finite odd order has positive Haar measure. Then $G$ has a prosolvable open normal subgroup.}} This motivates the following question:	{\sl{Let $G$ be a profinite group and suppose that the set of elements whose order, as a supernatural number, is odd has positive Haar measure. Is $G$ virtually prosolvable?}} We give a counterexample as follows.

\

Let $G=\prod_{n \geq 2}{\rm{SL}}(2,2^n)$ and denote by $\mathcal S_{\rm{odd}}$ the class of odd order finite solvable groups. Since $\mathcal S_{\rm{odd}}(G) \subseteq \mathcal S (G)$
and $R(G)=1,$ it follows from \cref{descs} that 
$\mathcal S_{\rm{odd}}(G) \subseteq \mathcal S(G)$ is contained in the $FC$-center of $G$, which coincides with the restricted direct product of the factors.
Moreover if $g=(g_n)_{n\in \mathbb N}\in \mathcal S_{\rm{odd}}(G)$, then $|g_n|$ is odd for every $n\in \mathbb N$ (otherwise $\mathcal S_{{\rm{odd}},_G}(g) =\varnothing).$ Thus 
$\mathcal{S}_{\rm{odd}}(G)$ is contained in the subset
$\Omega_{\rm{odd}}(G)$ consisting of the elements $(g_n)_{n\in \mathbb N}$
such that $|g_n|$ is odd for every $n\in \mathbb N$
and $g_n=1$ for all but finitely many $n\in \mathbb N.$ We claim that the equality $\mathcal{S}_{\rm{odd}}(G) = \Omega_{\rm{odd}}(G)$ holds. The group ${\rm{SL}}(2,2^n)$ has order $2^n(2^{2n}-1)$ and it contains precisely $2^{2n}-1$ nontrivial $2$-elements, %(see for example \cite{moller}) 
all of them have order $2$. Moreover, the order of an arbitrary element of ${\rm{SL}}(2,2^n)$ is either $2$ or odd. So the probability that an element of ${\rm{SL}}(2,2^n)$ has odd order is precisely $1-\frac{1}{2^n}.$
Then the probability that a randomly chosen element of $G$ has odd order (as a supernatural number) is
$$\eta:=\prod_{n \in \mathbb N}\left(1-\frac{1}{2^n}\right)>0.$$
It follows that the probability of randomly choosing an odd order element of $G$ is positive while $G$ is not virtually prosolvable.

\section{The class of finite nilpotent groups}\label{nilpot}

In this section we prove Theorems \ref{nilpp}, \ref{localnilp}, \ref{classp} and \ref{nofiniteorder}. If $G$ is a profinite group and $x\in G,$ we will denote by $\nil_G(x)$ the \emph{nilpotentizer} of $x$ in $G$, the set of elements $y \in G$ such that $\langle x,y \rangle$ is pronilpotent. Given a finite group $G$, a normal subgroup $N$ of $G$ and an element $x \in G$, the first thing we do is find a way to nicely bound $\mu_G(\mathcal N_G(x))$ from above in terms of $\mu_{G/N}(\mathcal N_{G/N}(xN))$.

\begin{lemma}\label{nilabcf}
	Let $G$ be a finite group, let $p$ be a prime and suppose that $N$ is an abelian minimal normal subgroup of $G$ whose order is a $p$-power. Let $x\in G$ and write
	$x=x_1x_2$ where $|x_1|$ is a $p$-power, $|x_2|$ is coprime with $p$ and $[x_1,x_2]=1.$
	Then
	$$\frac{|\nil_G(x)|}{|G|}\leq \frac{|\nil_{G/N}(xN)|}{|G/N|}\frac {|C_N(x_2)|}{|N|}.$$
\end{lemma}

\begin{proof}
   Given $y\in \nil_G(x)$ we want to estimate the size of $\Delta:=\nil_G(x)\cap Ny.$ We write $x=x_1x_2$ and $y=y_1y_2$, where $|x_1|$, $|y_1|$ are $p$-powers, $|x_2|, |y_2|$ are coprime with $p$ and $[x_1,x_2]=[y_1,y_2]=1.$ Since $\langle x,y
	\rangle$ is nilpotent, $[x_2,y_1]=[x_1,y_2]=1.$  Now assume $z\in \Delta.$ Then there exist $n_1, n_2\in N$ such that $z=y_1n_1y_2n_2,$ where $|y_1n_1|$ is a $p$-power, $|y_2n_2|$ is coprime with $p$ and $[y_1n_1,y_2n_2]=1.$
	First notice that $H = \langle x_2,y_2 \rangle$ and $K = \langle x_2, y_2n_2\rangle$ are Hall $p'$-subgroups of the solvable group $\langle x_2,y_2 \rangle N$, so by Hall's theorem there exists $n \in N$ such that $H^n = K$. Since $N \cap K = \{1\}$, if $h \in H$ then $|hN \cap K| \leq 1$, moreover $h^n \in hN$, so ${x_2}^n = x_2$ and ${y_2}^n = y_2n_2$. In other words, $y_2n_2 = {y_2}^{n}$ for some
	$n\in C_N(x_2).$ Thus the number of possible choices for
	$y_2n_2$ is at most $|C_N(x_2)|/|C_N(x_2)\cap C_N(y_2)|.$ Moreover
	$\langle x_2, y_1n_1\rangle$ nilpotent implies that $n_1\in C_N(x_2).$ Finally, using the commutator equalities $[ab,c]=[a,c]^b [b,c]$ and $[c,ab]=[c,b][c,a]^b$, since $N$ is abelian we have 
    \begin{align*}
        1 & = [y_1n_1,y_2n_2] = [y_1,n_2]^{n_1} [y_1,y_2]^{n_2n_1} [n_1,n_2] [n_1,y_2]^{n_2} = [y_1,n_2][n_1,y_2],
    \end{align*}
    thus $[y_1,n_2]=[y_2,n_1].$ Once $n_2$ is fixed, we count how many $n_1$ there are in $C_N(x_2)$ with $[y_1,n_2]=[y_2,n_1].$ It is possible that there is no $n_1$ satisfying this equation. In any case assume that 
	$[y_2,n_1]=[y_2,\tilde n_1]$, where $n_1$ and $\tilde n_1$ are two different elements of $C_N(x_2).$ Then $n_1\tilde n_1^{-1}\in C_N(y_2)\cap C_N(x_2)$, so, given $n_2$ there are at most $|C_N(x_2)\cap C_N(y_2)|$ possible choices for $n_1$. We conclude
	that the number of possibilities for the pair $(n_1,n_2)$
	is at most
	$$\frac{|C_N(x_2)|}{|C_N(x_2)\cap C_N(y_2)|}\cdot |C_N(x_2)\cap C_N(y_2)|=|C_N(x_2)|.
	$$
	Hence $|\Delta|\leq |C_N(x_2)|$. This implies the result.
\end{proof}

In the next proofs we will use the following notation. Given a profinite group $G$ and an element $g \in G$, for every prime $p$, $g=g_pg_{p^\prime}$ with $g_p$ a $p$-element, $g_{p^\prime}$ a $p^\prime$-element and $[g_p,g_{p^\prime}]=1.$ In particular, $g_p$ and $g_{p'}$ belong to $\langle g \rangle$. For every $g\in G$ and every open normal subgroup $N$ of $G$ we define the following non-negative integers:
\begin{enumerate}
	\item $\tau_{nonab}(g,N)$ is the number of nonabelian factors $X/Y$ in a chief series of $G/N$ such that $[X,g]\not\leq Y.$
	\item $\tau_{ab,p}(g,N)$ is the number of abelian factors $X/Y$ in a chief series of $G/N$ of order a $p$-power for some prime $p$ and such that $[X,g_{p^\prime}]\not\leq Y.$
	\item $\tau(g,N)=\tau_{nonab}(g,N)+\sum_p\tau_{ab,p}(g,N)$.
\end{enumerate}
	
	\begin{cor}\label{cencf}
		Let $G$ be a profinite group and let $g \in G$. If $\mu(\nil_G(g))>0,$ then there exists $u\in \mathbb N$ such that $\tau(g,N)\leq u$ for every open normal subgroup $N$ of $G.$
	\end{cor}
	
\begin{proof}
If $\mu(\nil_G(g))>0,$ then $\mu(\mathcal S_G(g))>0$ so by \cite[Corollary 1.3]{fgg} $g$ centralizes almost all the nonabelian chief factors of $G$, in other words there exists a constant $C$ such that $\tau_{nonab}(g,N) \leq C$ for every $N \unlhd_{\circ} G$. Moreover it follows from Lemma \ref{nilabcf} that, for every prime $p$, $g_{p^\prime}$ centralizes almost all the abelian chief factors of $G$ whose order is a power of $p$. Explicitly, if $N \unlhd_{\circ} G$ and $N = N_k \unlhd \cdots \unlhd  N_1 = G$ is a chief series of $G/N$, then, whenever $N_i/N_{i+1}$ is abelian, set $c_i := |C_{N_i/N_{i+1}}(g_{p^\prime})|/|N_i/N_{i+1}|$, being $p$ the unique prime divisor of $|N_i/N_{i+1}|.$
Note that, if $c_i < 1$, then $c_i \leq 1/2$, so $0 < \varepsilon = \mu(\nil_G(g)) \leq (1/2)^t$ where $t=\sum_p\tau_{ab,p}(g,N).$ It follows that $$\tau(g,N) = \tau_{nonab}(g,N) + t \leq C+\log_2(1/\varepsilon).\qedhere$$
\end{proof}

We need the following revised version of \cite[Theorem 25]{pel}.

\begin{thm}\label{thm_dh} Let $G$ be a finite group and $p$ a prime divisor of $G$. Then a $p$-element $g$ of $G$ centralizes all the abelian $p'$-chief factors and all the nonabelian chief factors of $G$ if and only if $g \in O_p(G).$
\end{thm}
\begin{proof}
Since the Fitting subgroup of $G$ is generated by the subgroups $O_p(G)$ where $p$ is a prime dividing $|G|$, the implication $(\Leftarrow)$ follows from \cite[Theorem 13.8(b)]{dh}. Suppose that a $p$-element $g$ of $G$ centralizes all the abelian $p'$-chief factors and all the nonabelian chief factors of $G$. This is equivalent to saying that $g$ centralizes all the chief factors of $G$ whose order is not a $p$-power.
	Let $q$ be a prime divisor of $|G|$ and let $C_q(G)$ be the intersection of the centralizers of the chief factors of $G$ whose order is divisible by $q
	.$ By \cite[Theorem 13.8(a)]{dh}, we have that $C_q(G)=O_{q^\prime,q}(G)$, where $O_q(G/O_{q^\prime}(G)) = O_{q^\prime,q}(G)/O_{q^\prime}(G)$. So it is sufficient to prove that a $p$-Sylow subgroup $P$ of $C=\cap_{q\neq p}O_{q^\prime,q}(G)$ is normal in $C$ and coincides with $O_p(G).$
	Since $O_{q^\prime,q}(G)$ is $q$-nilpotent, it follows from \cite[Theorem 13.4(a)]{dh}
	that $C$ is $q$-nilpotent for every $q\neq p.$ If $K_q$ is the normal $q$-complement in $C$, then $P=\cap_{q\neq p}K_q$.  Since $C$ is a normal subgroup of $G$, it follows that $P\leq O_p(G).$ On the other hand $O_p(G)\leq O_{q^\prime,q}(G)$ for every $q\neq p,$ so $O_p(G)\leq P.$
\end{proof}

So arguing as in \cite[Theorem 27]{pel}, we obtain the following.

% \begin{thm}
% 	Let $G$ be a profinite group and let $g$ be an element of $G$. %Suppose either that $p$ is odd or that $G$ is virtually prosolvable. 
% 	If $\mu(\nil_G(g))>0,$ then there is an open subgroup $H$ of $G$ such that  $g \in \fit(H).$ In particular, $g$ has finite order modulo $\fit(G).$ 
% \end{thm}
\begin{proof}[Proof of \cref{nilpp}]
It follows from \cref{cencf} that $G$ contains an open normal subgroup $M$ such that, for every prime $p$, $g_p$  centralizes all the chief factors of $M$ that are either nonabelian or have order coprime to $p.$ Hence, by \cref{thm_dh}, $g_p\in O_p(M\langle g\rangle)$. This is because, if $H$ is a closed subgroup of $G$, $P$ varies in the family of $p$-Sylow subgroups of $H$ and $N$ varies in the family of the open normal subgroups of $H$, then $O_p(H) = \cap_P P = \cap_P \cap_N PN = \cap_N \cap_P PN$ and $\cap_P PN$ is the preimage of $O_p(HN/N)$ in $H$. This implies that $g\in \fit(M\langle g\rangle).$
Moreover, $\langle g\rangle \cap M\leq \fit(M \langle g \rangle) \cap M = \fit(M)\leq \fit(G)$ so the order of $g$ modulo $\fit(G)$ is at most $|G/M|.$    
\end{proof}



In the proof of the next lemma we will use the following definition: if $K$ is a closed normal subgroup of $G$ and $p$ is a prime, we will denote by $\Delta_{G,p}(K)$ the set of the elements $g\in G$ with the following property: if $Y$ is an open normal subgroup of $K$, $X/Y$ is a chief factor of $K/Y$
and either $X/Y$ is nonabelian or has order coprime with $p$, then $[g,X]\leq Y.$

Recall that a profinite group is countably based if and only if it admits a countable family of open normal subgroups with trivial intersection.

\begin{lemma}\label{nilpel}
	Let $G$ be a countably based profinite group and let $P$ be a $p$-Sylow subgroup of $G.$  Suppose that $\mu(\nil_G(g))>0$ for every $g\in P.$
    Then the following hold:
    \begin{enumerate}
    \item $N_G(P)$ is an open subgroup of $G.$
    \item $P/O_p(G)$ is finite.
    \end{enumerate}
\end{lemma}
\begin{proof}
Since $G$ is countably based, there exists a descending chain  $\{N_i\}_{i\in \mathbb N}$ of open normal subgroups of $G$ such that $N_0=G,$  $\cap_{i\in \mathbb N}N_i=1$ and $N_i/N_{i+1}$ is a chief factor of $G/N_{i+1}$ for every $i\in \mathbb N.$    
Let $\Delta$ be the set of the natural numbers $u \in \mathbb{N}$ such that $N_{u}/N_{u+1}$ is a $p'$-chief factor or a nonabelian chief factor and let $\Lambda$ be the set of cofinite subsets of $\Delta$. For every $\lambda \in \Lambda$, we have that $C_\lambda := \cap_{i \in \lambda} C_P(N_i/N_{i+1})$ is a closed subgroup of $P$ and, by \cref{cencf}, it follows that $P= \cup_\lambda C_\lambda$. By the Baire Category Theorem, there exist $\lambda \in \Lambda$, an open normal subgroup $M$ of $P$ and $g \in P$ such that $gM \subseteq C_\lambda$. Since $C_\lambda$ is a subgroup of $G,$ this implies that $M$ itself is contained in $C_\lambda$. In particular, we can find an open normal subgroup $N$ of $G$ such that $M \leq \Delta_{G,p}(N)$ and therefore $M\cap N\leq O_p(N)\leq O_p(G).$ Hence $MO_p(G)/O_p(G)\cong M/(M\cap O_p(G))$ is an epimorphic image of 
$M/(M \cap N) \cong MN/N$ and therefore is finite, so $P/O_p(G)$ is finite too, which proves item (2). Write $P/O_p(G)=\{g_1O_p(G), \ldots, g_tO_p(G)\}$. By  \cref{cencf}, for each $1\leq i\leq t$, there exists an open normal subgroup $N_i$ of $G$ such that $g_i \in \Delta_{G,p}(N_i)$. Let $L = N \cap N_1 \cap \dots \cap N_t.$ It follows from \cref{thm_dh} that $P \leq O_p(PL)$ and, since $P$ is a Sylow subgroup of $G$, equality holds. Thus, $N_G(P)$ contains $L$ and consequently $N_G(P)$ is open in $G.$ This proves item (1).
\end{proof}

\begin{lemma}\label{intersecochiusi}Let $H$ and $K$ be closed subgroups of a profinite group $G$. If $H\cap K$ has infinite index in $H$ then there exists a closed normal subgroup $N$ of $G$ such that
\begin{enumerate}
\item $G/N$ is countably based;
\item $HN \cap KN$ has infinite index in $HN.$
\end{enumerate}
\end{lemma}
\begin{proof}
Let $\mathcal N$ be the family of the open normal subgroups of $G$. We have $H=\cap_{R\in \mathcal N}HR$ and $K=\cap_{S\in \mathcal N}KS,$
hence $$H \cap K=\cap_{R,S \in \mathcal N}HR\cap KS=\cap_{R,S\in \mathcal N}H(R\cap S) \cap K(R\cap S)=\cap_{T\in \mathcal N}HT\cap KT.$$
Since the index of $H\cap K$ in $H$ is infinite, there exists a countable subset $\{x_n\}_{\in \mathbb N}$ of $H$ such that $x_i(H\cap K) \neq x_j(H\cap K)$ whenever $i\neq j.$ In particular
for every $i\neq j,$ there exists an open normal subgroup $N_{i,j}$ of $G$ such that $x_i^{-1}x_j\notin HN_{i,j}\cap KN_{i,j}.$ Let $N=\cap_{i,j}N_{i,j}.$ By construction, $x_i^{-1}x_j\notin HN\cap KN$ whenever $i\neq j$.
\end{proof}

% \begin{thm}
% 	Let $G$ be a  profinite group and let $P$ be a $p$-Sylow subgroup of $G.$  Suppose that $\mu(\nil_G(g))>0$ for every $g\in P.$
% 	Then the following hold:
% 	\begin{enumerate}
% 		\item $N_G(P)$ is an open subgroup of $G.$
% 		\item $P/O_p(G)$ is finite.
% 		\item $O_p(G)C_G(g)$ is open in $G$ for every $g \in P$.
% 		\item If $P$ is finite, then $C_G(P)$ is open in $G$.
% 	\end{enumerate}
% \end{thm}

\begin{proof}[Proof of \cref{localnilp}]
	Assume by contradiction that $H=N_G(P)$ is not open in $G.$ Then $H$ has infinite index in $G$, and therefore there exists a countable subset $\{g_n\}_{n\in \mathbb N}$ of $G$ such that $g_iH \neq g_jH$ whenever $i\neq j.$ Since $H$ is closed in $G$, for every pair $(i,j)$ with $i\neq j$ there exists an open normal subgroup $N_{i,j}$ of $G$ with $g_i^{-1}g_j\notin HN_{i,j}.$ Let $N=\cap_{i,j}N_{i,j}.$ Then $G/N$ is countably based, so by \cref{nilpel} $N_G(PN)$ has finite index in $G$. However $g_i^{-1}g_j$ does not normalize $PN$ if $i \neq j$, because $N_G(PN)=N_G(P)N=HN \subseteq HN_{i,j}$ where the first equality follows from the Frattini argument. This contradiction proves item (1). Now let $Q$ be another $p$-Sylow subgroup of $G$. Assume, by contradiction, that $|P:P\cap Q|$ is not finite. Then, by \cref{intersecochiusi}, there exists a closed normal subgroup $N$ of $G$ that $G/N$ is countably based and $|PN/N:PN/N\cap QN/N|=\infty,$ against \cref{nilpel} (2). Hence $|P: P\cap P^g|$ is finite for every $g\in G$. Since, by (1), $P$ has only finitely many conjugates in $G$, we deduce that $P/O_p(G)$ is finite, proving (2). Let $\Omega_p(G)$ be the set of $p$-elements of $G$. Since $|P/O_p(G)|$ and $|G:N_G(P)|$ are finite, there exists a finite subset $T$ of $G$ such that $\Omega_p(G) = \cup_{t\in T} tO_p(G).$ Now let $g\in \Omega_p(G).$ Since the Haar measure is invariant under conjugation, $\mu(\mathcal{N}_G(g))>0$. If $x\in \nil_G(g),$ then $x=x_1x_2$, where $x_1\in \Omega_p(G)$ and $[x_2,g]=1.$ Hence
    $$\nil_G(g)\subseteq \Omega_p(G)C_G(g)=\bigcup_{t\in T}tO_p(G)C_G(g).$$
	But then 
    \begin{align*}
    \eta & =\mu(\nil_G(g))\leq \mu(\cup_{t\in T}tO_p(G)C_G(g))=|T|\mu(O_p(G)C_G(g)) =\frac{|T|}{|G:O_p(G)C_G(g)|}
    \end{align*}
    and since $\eta > 0$, we conclude that 
	$|G:O_p(G)C_G(g)|^{-1}\geq \eta/|T|.$ This proves item (3). If $P$ is finite, then $O_p(G)$ is finite too hence $C_G(g)$ is open in $O_p(G) C_G(g)$ for every $g \in P$. Now item (3) implies that $C_G(g)$ is open in $G$ and, since $P$ is finite, $C_G(P) = \cap_{g \in P} C_G(g)$ is open too. This proves item (4).
\end{proof}

% \begin{prop}For every positive real number $\varepsilon$
% there exists a solvable profinite group $G$ which contains an element $g$ such that $\mu(\nil_G(g))>1-\varepsilon$,
% but $g$ has infinite order modulo $Z_\infty(G).$
% In particular, there is no open subgroup $H$ with $g\in Z_\infty(H)$.
% \end{prop}

Next, we show that \cref{nilpp} can't be generalized by considering the hypercenter instead of the Fitting subgroup.

\begin{proof}[Proof of \cref{nofiniteorder}]
Let $p$ and $q$ be two primes with $p$ dividing $q-1$ and consider the group $Y_t$ (associated to the primes $p$ and $q$) defined in the proof of \cref{allfingrps} (with the same notation) so that $|Y_t|=q^{n_t} p^{t+1}$ where $n_t=p^t$.  All the elements of $Y_t$ that are not contained in $N$ have order a $p$-power, and therefore $$\cfrac{|\Omega_p(Y_t)|}{|Y_t|} \geq \cfrac{|Y_t|-|N|}{|Y_t|} = 1-\cfrac{1}{p^{t+1}}.$$
 A $p$-Sylow subgroup $R$ of $Y_t$ has trivial normal core and $|Y_t:R|=q^{n_t}$. Thus, $Y_t$ embeds as a transitive subgroup into $\sym(q^{n_t})$ and we may consider the group $X_t = W_t \rtimes Y_t < C_{p^t} \wr Y_t$ where $C_{p^t}$ is cyclic of order $p^t$ and $W_t = \{(x_1,\ldots,x_{q^{n_t}}) \in (C_{p^t})^{q^{n_t}}\ |\ \prod_{i=1}^{q^{n_t}}x_i=1\}$. Then $X_t$ has trivial center, since the transitivity (and the faithfulness) of $Y_t$ implies that any element of $Z(X_t)$ has the form $(a,\ldots,a) \in (C_{p^t})^{q^{n_t}}$, and $a^{q^{n_t}} = \prod_{i=1}^{q^{n_t}}a = 1$ implies that $a=1$. Hence $\Z_\infty(X_t)=1.$ If we take any $g \in W_t \smallsetminus \{1\}$ and any $y \in \Omega_p(X_t)$, then $\la g,y\ra \leq \la O_p(X_t),y\ra$, which is a $p$-group, so $y \in \mathcal N_{X_t}(g)$. Notice also that every element in $W_t \Omega_p(Y_t)$ is a $p$-element and that
\[W_t \Omega_p(Y_t) = \bigcup_{b \in W_t} b\Omega_p(Y_t)\]
is a disjoint union. Therefore,
\[\cfrac{|\nil_{X_t}(g)|}{|X_t|} \geq \cfrac{|W_t||\Omega_p(Y_t)|}{|W_t||Y_t|}  \geq 1-\cfrac{1}{p^{t+1}}.\]
Now choose $s\in \mathbb N$ such that
$$\prod_{t \geq s}\left( 1- \cfrac{1}{p^{t+1}} \right) > 1-\varepsilon$$ and take the profinite group $G=\prod_{t\geq s} X_t.$ 
Clearly $Z_\infty(G)=\prod_{t\geq s} Z_\infty(X_t)=1$.
Now consider an element $g = (g_t)_{t \geq s}$ such that $g_t \in W_t$ has order $p^t$ for all $t$. Clearly $g$ has infinite order. On the other hand, $\nil_G(g)=\prod_{t\geq s}  \nil_G(g_t)$ and therefore
$$\mu(\nil_G(g)) \geq \prod_{t \geq s}\left( 1- \cfrac{1}{p^{t+1}} \right) > 1-\varepsilon.$$
For every $t\in \mathbb N,$ $X_t$ is solvable with derived length 4, so $G$ is solvable.

If $H$ is any open subgroup of $G$, then $g \not \in Z_{\infty}(H)$. Indeed, assume this is not true. There exists $r \geq s$ such that $N = \prod_{t \geq r} X_t \leq H$ and $N$ is an open normal subgroup of $G$. Now $\langle g \rangle \cap N \leq Z_{\infty}(H) \cap N \leq Z_{\infty}(N) = 1$, a contradiction since $g$ has infinite order.
\end{proof}

Finally, we prove a generalization of \cite[Theorem 29]{pel}.

\begin{proof}[Proof of \cref{classp}]Suppose $\mu(\mathcal P_p(G))>0$. Since $\mathcal P_p(G)\subseteq \mathcal S(G),$ it follows from \cref{descs} that $G$  is virtually prosolvable.
Let $C$ be the set of the elements $g$ in $G$ which centralize almost all the abelian chief factors of $G$ of order coprime to $p$, that is, there exists a positive integer $k$ such that, for every $N \unlhd_{\circ} G$, the set of abelian $p'$-chief factors of $G/N$ not centralized by $g$ has cardinality at most $k$. Then $C$ is a measurable subgroup of $G$. It follows from \cite[Corollary 24]{pel} that $\mathcal P_p(G)\subseteq C$, so $C$ has finite index and is closed in $G$. We claim that $C$ is virtually pro-$p.$ It is not restrictive to assume that $C$ is prosolvable and $O_p(C)=1.$ First, repeating the  arguments in \cite[Theorems 29 and 30]{pel} it can be proved that a $p$-Sylow subgroup $P$ of $C$ is finite. Since $P\leq C,$ every element of $P$ centralizes almost all the $p^\prime$-chief factors of $C$ and therefore there exists an open  normal subgroup $D$ of $C$ such that $P$ centralizes all  the $p^\prime$-chief factors of $D$. By \cite[Corollary 26]{pel} $P = O_p(PD)$, and therefore $D \subseteq N_C(P)$ so $|C:N_C(P)|$ is finite. But then $C$ contains only finitely many $p$-elements, so $\mathcal P_{p,C}(g)$
is finite for any $g \in C$. Thus, from the assumption $\mu(\mathcal P_p(G))>0$, we conclude that $C$, and hence $G$, is finite.
\end{proof}

%\textcolor{red}{Mi sembra che adesso le cose scritte nel paragrafo successivo sono già scritte nell'introduzione e quindi il paragrafo possa/debba essere tolto}

%We make a final remark. In \cite{aemp}, it is observed that the condition $\mathcal N(G) = G$ is stronger than $G$ being virtually pronilpotent; indeed, it suffices to consider the example $\mathbb Z_p \rtimes C_2 = \mathbb Z_p \rtimes \la x \ra$, where $p$ is an odd prime and $x$ acts on the $p$-adic integers by inversion, and observe that $\mathcal N_G(x) = \{1,x\}$. So it is natural to ask if $\mu(\mathcal N(G)) > 0$ if, and only if, $G$ is virtually pronilpotent. Compared to \cref{descs}, this is considerably harder to approach because we do not have a group-theoretical characterization of the {nilpotentizer} like we do for the {solvabilizer} nor do we know if $\mathcal N(G)$ is a subgroup. Hence, it would be desirable to find answers to \cref{q1} and \cref{q3} for the class of finite nilpotent groups.

\section{Measurable sets related to pronilpotency} \label{Lambda}

The proof of \cref{nobaer} relies on the following remark:

\begin{lemma}\label{sumcent}Let $G$ be a finite group, $p$ a prime, and let $N$ be a normal subgroup of $G$ whose order is coprime with $p.$ If $x$ is a $p$-element of $G$ and $y\in \Lambda_{G,p}(x)$, then
$\Lambda_{G,p}(x) \cap yN=yC_N(x^y)C_N(x)$.
\end{lemma}

\begin{proof}
Notice that $P:=\langle x, x^y\rangle$ is a $p$-Sylow subgroup of $PN.$ Now if $yn \in \Lambda_{G,p}(x) \cap yN,$ then $\langle x, x^{yn}\rangle=P^m$ for some $m\in N.$ If $h \in P$, then $|hN \cap P^m| \leq 1$ and moreover $h^m \in hN$, so $x^m=x$ and $x^{ym}=x^{yn}$. Therefore $m \in C_N(x)$ and
$mn^{-1}\in C_N(x^y).$ It follows that $yn \in \Lambda_{G,p}(x) \cap yN$ if and only if $n\in C_N(x^y)C_N(x)$.
\end{proof}

%\begin{proof}[Proof of \cref {nobaer}]
%Let $q$ be a prime such that $p$ divides $q-1$ and choose $n\in \mathbb N$ such that $n-1\geq \eta n.$
%${qn-q+1}\geq \eta qn.$ 
%Consider a cyclic
%group $A=\langle a\rangle$ of order $p$ and let $\mathbb F_q$ be the field with $q$ elements. Choose $\lambda \in \mathbb F_q^\times$ with $|\lambda|=p.$ Then
%$A$ acts on $\mathbb F_q$ by setting $f^a=\lambda f$ for every $f\in \mathbb F_q.$ Now let $\sigma=(1,2,\dots,n)\in \perm(n)$ and let $C=\langle \sigma \rangle \leq \perm(n).$ The action of $A$ on $\mathbb F_q$ gives rise to an action of the wreath product $H=A\wr C$ on the vector space $V=\mathbb F_q^n$ of dimension $n$ over $\mathbb F_q.$ We consider the semidirect product $G=V\rtimes H \leq {\rm{AGL}}(n,q).$ Since $\fit(G)=V,$ we have $O_p(G)=1.$ Let $x=(a,1,\dots,1)\in A^n\leq H.$
%Then $$C_V(x)=\{(0,f_2,\dots,f_n)\mid f_2,\dots,f_n\in V\}.$$
%For $1\leq i\leq n-1$ let $\Sigma_i:=A^n\sigma^i\ \subseteq H.$ If $g\in \Sigma_i$, we have
%$$C_V(x^g)=C_V(x^{\sigma^i})=\{(f_1,\dots,f_i,0,f_{i+2},\dots f_n)\mid f_1,\dots,f_i,f_{i+2},\dots,f_n\in V\}.$$
%In particular, $C_V(x)+C_V(x^{g})=V$. Since $g\in \Lambda_{G,p}(x)$, it follows from \cref{sumcent} that $Vg\subseteq \Lambda_{G,p}(x).$
%We have so proved that
%$\bigcup_{1\leq i\leq n-1}VA^n\sigma^i \subseteq \Lambda_{G,p}(x).$ 
%Moreover $C_G(x)\cong \mathbb F_q^{n-1}\langle x\rangle \subseteq \Lambda_{G,p}(x).$
%But then
%$$\frac{|\Lambda_{G,p}(x)|}{|G|}\geq\frac{n-1}{n}\geq \eta. \qedhere$$
%+\frac{1}{qn}=\frac{qn-q+1}{qn}\geq \eta.\qedhere $$
%\end{proof}

\begin{proof}[Proof of \cref{nobaer}]
Choose a prime $q$ such that $p^t$ divides $q-1$ and choose a positive integer $n$ such that $n-1\geq\eta n$. Let $A=\la a \ra$ be the cyclic group of order $p^t$,  let $\mathbb F_q$ be the field with $q$ elements and choose $\lambda \in \mathbb F_q^\times$ of order $p^t$. Then $A$ acts on $\mathbb F_q$ by $f^a = \lambda f$ for every $f \in \mathbb F_q$. Now, let $\sigma=(1,2,\dots,n)\in \perm(n)$
and let $C=\langle \sigma \rangle \leq \perm(n).$ The action of $A$ on $\mathbb F_q$ gives rise to an action of the wreath product $H=A\wr C$ on the vector space $V=\mathbb F_q^n$ of dimension $n$ over $\mathbb F_q.$ We consider the semidirect product $G=V\rtimes H \leq {\rm{AGL}}(n,q).$ {Since $H$ acts faithfully on $V$, it follows that $O_p(G)=1$.} Let $x=(a,1,\dots,1)\in A^n\leq H.$
Then $$C_V(x)=\{(0,f_2,\dots,f_n)\mid f_2,\dots,f_n\in V\}.$$
For $1\leq i\leq n-1$
let $\Sigma_i:=A^n\sigma^i\ \subseteq H.$ If $g\in \Sigma_i$, we have
$$C_V(x^g)=C_V(x^{\sigma^i})=\{(f_1,\dots,f_i,0,f_{i+2},\dots f_n)\mid f_1,\dots,f_i,f_{i+2},\dots,f_n\in V\}.$$
In particular, $C_V(x)+C_V(x^{g})=V$. Since $g\in \Lambda_{G,p}(x)$, it follows from \cref{sumcent} that $Vg\subseteq \Lambda_{G,p}(x).$
We have so proved that
$\bigcup_{1\leq i\leq n-1}VA^n\sigma^i \subseteq \Lambda_{G,p}(x).$ 
%Moreover $C_G(x)\cong \mathbb F_q^{n-1}\langle x\rangle \subseteq \Lambda_{G,p}(x).$
But then
$$\frac{|\Lambda_{G,p}(x)|}{|G|}\geq\frac{n-1}{n}\geq \eta. \qedhere$$
\end{proof}

As we have already recalled, by \cite[Theorem 1]{aemp} if $\mu(\nil_G(g))>0$ for every $g\in G$ then $G$ is virtually pronilpotent. This motivates the following question:
\begin{question}
	Let $G$ be a profinite group. Suppose that $\mu(\Lambda_{G}(g))>0$ for every  $g\in G.$ Is $G$ virtually pronilpotent?
\end{question}

We also propose the following related question.

\begin{question}
    Let $G$ be a profinite group and suppose that $$\mu(\{(x,g) \in G^2 \mid \la x,x^g\ra \text{ is pronilpotent}\})>0.$$ Is $G$ virtually pronilpotent?
\end{question}


\begin{thebibliography}{1}
	\bibitem{az} A. Abdollahi and M. Zarrin, {Non-nilpotent graph of a group}, Comm. Algebra 38 (2010) 4390--4403.
%\bibitem{aschbacher} M. Aschbacher, Finite group theory, Cambridge Studies in Advanced Mathematics, 10, Cambridge University Press, Cambridge, 1986. 
%\bibitem{bhrd} J. Bray, D. Holt, C. Roney-Dougal, The maximal subgroups of the low-dimensional finite classical groups.With a foreword by Martin Liebeck. London Mathematical Society Lecture Note Series, 407. Cambridge University Press, Cambridge, 2013. 
	\bibitem{bm} W. R. Brian and M. W. Mislove,
	Every infinite compact group can have a non-measurable subgroup,
	Topology Appl. 210 (2016), 144–146.
%    \bibitem{bln} T. C. Burness, A. Lucchini and D. Nemmi, On the soluble graph of a finite group, Journal of Combinatorial Theory, Series A, Volume 194, 2023.
%    \bibitem{cc1} C. Casolo, Finite groups in which subnormalizers are subgroups, Rend. Sem. Mat. Univ. Padova 82 (1989), 25--53. 
%	\bibitem{cc2}  C. Casolo,  Subnormalizers in finite groups, Comm. Algebra 18 (1990), no. 11, 3791–3818.
%    \bibitem{engel} E. Detomi, A. Lucchini and D. Nemmi, The Engel graph of a finite group, Forum Math. 35 (2023), no. 1, 111–122.
	\bibitem{aemp} E. Detomi, A. Lucchini, M. Morigi and   P. Shumyatsky,
	Probabilistic properties of profinite groups, arXiv:2304.04573, to appear in Israel J. Math.
	\bibitem{dh} K.	Doerk and T. Hawkes, Finite soluble groups, De Gruyter Expositions in Mathematics, 4. Walter de Gruyter \& Co., Berlin, 1992. 
	\bibitem{fgg}J. Fulman, D. Garzoni and R. Guralnick, Probabilistic Generation of Finite Almost Simple Groups, 
Adv. Math. 452 (2024), Paper No. 109796, 33 pp.
	%\bibitem{ghe} P. Gheri,  On the number of $p$-elements in a finite group, Ann. Mat. Pura Appl. (4) 200 (2021), no. 3, 1231--1243.
% \bibitem{pen}   R. Guralnick, T. Penttila, C. Praeger and J. Saxl, Linear groups with orders having certain large prime divisors,Proc. London Math. Soc. (3) 78 (1999), no. 1, 167–214.
   \bibitem{gor} D. Gorenstein, Finite Simple Groups: An Introduction to Their Classification, University Series in Mathematics, Plenum Publishing Corp., New York, 1982.
\bibitem{gk} R. M. Guralnick and W. M. Kantor, Probabilistic generation of finite simple groups,
J. Algebra 234 (2000), no. 2, 743–792.
% \bibitem{huppert} B.   Huppert,  Endliche Gruppen, I. Die Grundlehren der mathematischen Wissenschaften, Band 134. Springer-Verlag, Berlin-New York, 1967. xii+793 pp
% \bibitem{jones}   G. Jones, Primitive permutation groups containing a cycle,  Bull. Aust. Math. Soc. 89 (2014), no. 1, 159–165.
%\bibitem{kl} P. Kleidman and M. Liebeck, The subgroup structure of the finite classical groups, London Mathematical Society Lecture Note Series, 129. Cambridge University Press, Cambridge, 1990.
	\bibitem{ls} M. Larsen and A. Shalev,  Hausdorff dimension and randomly free groups,
	Math. Ann. 371 (2018), no. 3-4, 1409--1427.
%    \bibitem{lss} J. C. Lennox and S. E. Stonehewer, Subnormal subgroups of groups, New York: Oxford University Press, 1987.
 %   \bibitem{lps} M. Liebeck, C. Praeger and J.  Saxl, The maximal factorizations of the finite simple groups and their automorphism groups. Mem. Amer. Math. Soc. 86 (1990), no. 432, 
 \bibitem{LP} L. L\'evai and L. Pyber, Profinite groups with many commuting pairs or involutions, Arch. Math. 75 (2000), 1--7. Birkh\"auser Verlag, Basel, 2000.
	\bibitem{solv} A. Lucchini,  Solubilizers in profinite groups, J. Algebra 647 (2024), 619–632.
	\bibitem{pel}  A. Lucchini and N. Otmen, $p$-elements in profinite groups,  Monatsh. Math. 206 (2025), no. 4, 933–948.
%	\bibitem{moller} J. M. M\o ller,  The number of $p$-elements in finite groups of Lie type of characteristic $p$, Proc. Amer. Math. Soc. 149 (2021), no. 7, 2805–2812.
%\bibitem{mp} G. L. Morgan, C. W. Parker, The diameter of the commuting graph of a finite group with trivial centre, {J. Algebra}{393} (2013), 41--59.
%    \bibitem{rib} L. Ribes, Profinite graphs and groups, Ergebnisse der Mathematik und ihrer Grenzgebiete. 3. Folge. A Series of Modern Surveys in Mathematics, 66, Springer, Cham, 2017.
%    \bibitem{rz} L. Ribes and P. Zalesskii, Profinite groups, Ergebnisse der Mathematik und ihrer Grenzgebiete. 3. Folge. A    Series of Modern Surveys in Mathematics 40, Springer-Verlag, Berlin, 2000.
 %   \bibitem{dsss} M. du Sautoy, D. Segal and A. Shalev (editors), New horizons in pro-$p$ groups, Birhkhauser Boston, Inc., Boston, 2000.
%\bibitem{pop}   F. Pop, On prosolvable subgroups of profinite free products and some applications, Manuscripta Math. 86 (1995), no. 1, 125–135. 
   % \bibitem{pass} D. Passman, Permutation groups. Revised reprint of the 1968 original. Dover Publications, Inc., Mineola, NY, 2012.
 %   \bibitem{reid} C. D. Reid, The Generalised Pro-Fitting Subgroup of a Profinite Group. Communications in Algebra, 41(1), 2013, 294–308.
  %  \bibitem{rob} D. J. S. Robinson, A course in the theory of groups, Graduate Texts in Mathematics, 80, Springer-Verlag New York, 1996.
    \bibitem{hyper} P.	Schmid, 
	{The hypercenter of a profinite group},
	Beitr. Algebra Geom. 55 (2014), 645--648.
	\bibitem{sh} A. Shalev,  Profinite groups with restricted centralizers, Proc. Amer. Math. Soc. 122 (1994), no. 4, 1279–1284.
    \bibitem{mv} M. Vannacci, On hereditarily just infinite profinite groups obtained via iterated wreath products, J. Group Theory 19 (2016), no.~2, 233--238.
  
	
%	\bibitem{wilson} J. S. Wilson,   Profinite Groups. Clarendon Press, Oxford, 1998.
%	\bibitem{zel} E.I. Zelmanov, On periodic compact groups, Israel J. Math. 77 (1992), 83--95.
\end{thebibliography}	
	
	\end{document}
